%% kapa: exam=1 round=20 year=2009 period=1 form=A subject=civil_law subject_ko=민법 seq=15 number=15
%% source: Q-Net 감정평가사 기출문제 (cst00309) artlSeq=5024178 file=extracted/1008/5024178/2009년 제20회 감정평가사 1차시험 기출문제/1교시.hwp
%% conversion: kordoc 4.24.51 -> tools/gen1.py (auto); review=proofread-sonnet; flags=none
\begin{question}{15}
甲 사단법인의 이사 乙은 丙으로부터 甲의 사업운영자금 명목으로 甲명의로 1억원을 차용한 후, 사적인 용도로 사용하였다. 다음 중 옳지 \uline{않은} 것은? (다툼이 있으면 판례에 의함)
\begin{choices}
  \item 乙이 다른 이사들의 동의 없이 단독으로 금전을 차용하였다고 해서 그 행위가 무효인 것은 아니다.
  \item 丙이 乙의 배임적 의도를 알 수 없었던 경우에는 乙의 대표행위는 甲의 행위로서 유효하므로 甲은 丙에게 금전소비대차에 기한 의무를 부담한다.
  \item 丙이 乙의 배임적 의도를 알고 있었던 경우에는 丙은 甲에게 금전소비대차에 기한 채무이행을 청구하지 못한다.
  \item 법인의 불법행위가 성립하는 경우 乙은 甲과 연대하여 丙에게 배상책임을 진다.
  \item 乙의 행위는 직무에 관한 행위이므로 丙이 乙의 배임적 의도를 알고 있었던 경우에도 丙은 甲의 불법행위책임을 물을 수 있다.
\end{choices}
\end{question}
\answer{5}
